% !TEX TS-program = lualatex
\documentclass[tikz,border=0.4mm]{standalone}
\usepackage{tikz}
\usetikzlibrary{shadows.blur}

\usepackage{gaplogo}

\def\NoTextMode{}

\begin{document}
\begin{tikzpicture}

% === Parameters ===
\def\discradius{16mm}        % disc radius

% shadow style
\tikzset{discShadow/.style={blur shadow={
            shadow blur steps=10,
            shadow xshift=0pt,
            shadow yshift=-0.7pt,
            shadow blur radius=1.2pt,
            shadow opacity=40}}}

\tikzset{logoShadow/.style={blur shadow={
            shadow blur steps=10,
            shadow xshift=0pt,
            shadow yshift=-1pt,
            shadow blur radius=1.2pt,
            shadow opacity=40}}}

\definecolor{outerGradientA}{gray}{0.8}
\definecolor{outerGradientB}{gray}{0.999}

\definecolor{innerGradientA}{gray}{0.8}
\definecolor{innerGradientB}{gray}{1.0}

% Background disc with linear gradient, and blurred shadow
\shade[shading=axis,
       left color=outerGradientA, right color=outerGradientA, % both sides same
       top color=outerGradientA, bottom color=outerGradientB,
       shading angle=180,
       discShadow]
  (0,0) circle (\discradius);

% Small disc with reverse linear gradient, to give the impression of a border
\shade[shading=axis,
       left color=innerGradientB, right color=innerGradientB,
       top color=innerGradientB, bottom color=innerGradientA,
       shading angle=180]
  (0,0) circle (\discradius*0.92);

% TODO: Render shadows for the GAP logo
\begin{scope}[yshift=-0mm]
  \coordinate (A) at (    \rotationangle:\radius) {};
  \coordinate (B) at (120+\rotationangle:\radius) {};
  \coordinate (C) at (240+\rotationangle:\radius) {};
  \coordinate (Z) at (0,0) {};

  % node shadows
  \path[logoShadow] (Z) circle[radius=\noderadius+1.1pt];
  \path[logoShadow] (A) circle[radius=\noderadius+1.1pt];
  \path[logoShadow] (B) circle[radius=\noderadius+1.1pt];
  \path[logoShadow] (C) circle[radius=\noderadius+1.1pt];

  % edge and arrow shadows
  \foreach \angle in {0,120,240} {
    % edge
    \path[logoShadow,rotate around={\angle+\rotationangle:(Z)}]
        (Z) ++(0,-\edgelinewidth/2) rectangle ++(\radius,\edgelinewidth);

    % double-arrows
    % first some setup...
    \tikzmath{
       \startangle = \angle+\rotationangle+\margin+17;
       \endangle = \angle+120+\rotationangle-\margin-17;
       \width = \arrowlinewidth/2;
    }
    % ... now draw the line part
    \path[logoShadow]
        (\startangle:\radiusB+\width) arc (\startangle:\endangle:\radiusB+\width) --
        (\endangle:\radiusB-\width) arc (\endangle:\startangle:\radiusB-\width) --
        cycle;

    % crude approximation of the first arrow
    \path[logoShadow]
        (\startangle-0.9:\radiusB)
        -- (\startangle+3:\radiusB+2.6*\width)
        to [bend left=60] (\startangle+0.5:\radiusB+3.4*\width)
        to [bend left=7] (\startangle-16.3:\radiusB+0.5*\width)
        to [bend left] (\startangle-16.3:\radiusB-0.5*\width)
        to [bend right=7] (\startangle:\radiusB-3.35*\width)
        to [bend left=90] (\startangle+3.5:\radiusB-2.6*\width)
        to [bend left=5] cycle
        ;

    % crude approximation of the second arrow
    \path[logoShadow]
        (\endangle+0.9:\radiusB)
        -- (\endangle-3:\radiusB+2.6*\width)
        to [bend right=60] (\endangle-0.5:\radiusB+3.4*\width)
        to [bend right=7] (\endangle+16.3:\radiusB+0.5*\width)
        to [bend right] (\endangle+16.3:\radiusB-0.5*\width)
        to [bend left=7] (\endangle:\radiusB-3.35*\width)
        to [bend right=90] (\endangle-3.5:\radiusB-2.6*\width)
        to [bend right=5] cycle
        ;
  }


\end{scope}

% Render the GAP logo atop
\begin{scope}
  % set coordinates of the four circles, starting at the one on the right
  \coordinate (A) at (    \rotationangle:\radius) {};
\ifx\mirrored\undefined
  \coordinate (B) at (120+\rotationangle:\radius) {};
  \coordinate (C) at (240+\rotationangle:\radius) {};
\else
  \coordinate (B) at (240+\rotationangle:\radius) {};
  \coordinate (C) at (120+\rotationangle:\radius) {};
\fi
  \coordinate (Z) at (0,0) {};

  % draw the edges connecting the four circles / nodes
  \draw[shorten >= 2mm, shorten <=2mm, dynkinedgeA] (Z) -- (A);
  \draw[shorten >= 2mm, shorten <=2mm, dynkinedgeB] (Z) -- (B);
  \draw[shorten >= 2mm, shorten <=2mm, dynkinedgeC] (Z) -- (C);

  % draw the nodes
  \path[dynkinnodeZ] (Z) circle[radius=\noderadius];
  \path[dynkinnodeA] (A) circle[radius=\noderadius];
  \path[dynkinnodeB] (B) circle[radius=\noderadius];
  \path[dynkinnodeC] (C) circle[radius=\noderadius];


  % draw the three double-ended arrows which indicate swapping of the circles
  % for this we use circle arcs but with slightly reduced radius
\draw[<->, dynkinarrowAB] ({\rotationangle+\margin}:\radiusB)
  arc ({\rotationangle+\margin}:{120+\rotationangle-\margin}:\radiusB);

\draw[<->, dynkinarrowBC] ({120+\rotationangle+\margin}:\radiusB)
  arc ({120+\rotationangle+\margin}:{240+\rotationangle-\margin}:\radiusB);

\draw[<->, dynkinarrowCA] ({240+\rotationangle+\margin}:\radiusB)
  arc ({240+\rotationangle+\margin}:{360+\rotationangle-\margin}:\radiusB);


  % Invisible circle to make bounding box centered on (0, 0)
  \draw[draw=none, fill=none] let
    \p1 = ($(current bounding box.east)$), \n1 = {veclen(\x1,\y1)},
    \p2 = ($(current bounding box.west)$), \n2 = {veclen(\x2,\y2)},
    \p3 = ($(current bounding box.south)$), \n3 = {veclen(\x3,\y3)},
    \p4 = ($(current bounding box.north)$), \n4 = {veclen(\x4,\y4)},
    \n5={max(\n1, \n2, \n3, \n4)}
    in (0, 0) circle (\n5);


\ifx\NoTextMode\undefined
    % draw the GAP text
    \node[anchor=west] at (1.2, 0.8) {\gapstyle{\textcolor{Gcolor}{G}}\extrastyle{roups}};
    \node[anchor=west] at (1.2, 0  ) {\gapstyle{\textcolor{Acolor}{A}}\extrastyle{lgorithms}};
    \node[anchor=west] at (1.2,-0.8) {\gapstyle{\textcolor{Pcolor}{P}}\extrastyle{rogramming}};
\fi

\end{scope}

\end{tikzpicture}
\end{document}
